%------------------------------------------------------------------------------%

\usepackage{calculo_varias_variaveis-1}

\title{Derivação Implícita}

%------------------------------------------------------------------------------%
\begin{document}

\addtitle

%------------------------------------------------------------------------------%
\section{Derivação Implícita}

%------------------------------------------------------------------------------%
\begin{frame}{Derivação Implícita}
  Suponha que $F(x, y)$ seja diferenciável e que a equação
  \[
    F(x, y) = 0
  \]
  defina $y$ como função de $x$, encontre \(\frac{dy}{dx}\)
\end{frame}

%------------------------------------------------------------------------------%
\begin{frame}{Derivação Implícita -- Justificativa}
\begin{columns}
\column{0.4\linewidth}
\begin{align*}
  \onslide<+->{w(x)          = F\big(x, y(x)\big)                              & = 0 \\[4mm]}
  \onslide<+->{\frac{dw}{dx} = \D{F}{x} \frac{dx}{dx} + \D{F}{y} \frac{dy}{dx} & = 0}
\end{align*}
\column{0.6\linewidth}
\begin{align*}
  \onslide<+->{F_x \frac{dx}{dx} + F_y \frac{dy}{dx} & = 0      \\[2mm]}
  \onslide<+->{F_x + F_y \frac{dy}{dx}               & = 0      \\[2mm]}
  \onslide<+->{F_y \frac{dy}{dx}                     & = - F_x  \\[2mm]}
  \onslide<+->{\frac{dy}{dx}                         & = - \frac{F_x}{F_y} 
               \qquad \text{desde que } F_y \neq 0}
\end{align*}
\end{columns}
\end{frame}

%------------------------------------------------------------------------------%
\begin{frame}{Derivação Implícita -- Fórmula}
  Suponha que $F(x, y)$ seja diferenciável e que a equação
  \[
    F(x, y)=0
  \]
  defina $y$ como função de $x$, \(y=y(x)\), então
  \[
    \frac{dy}{dx} = -\frac{F_x}{F_y}
  \]
  desde que \(F_y \neq 0\)
\end{frame}

%------------------------------------------------------------------------------%
\addtocounter{exemplo}{1}
\begin{frame}{Exemplo \theexemplo}
  Calcule \(\frac{dy}{dx}\) se \(y^2-x^2-\sin(xy) = 0\)
\end{frame}

\begin{frame}{Exemplo \theexemplo}
  \[
    \frac{dy}{dx} = -\frac{F_x}{F_y}
  \]
  \pause
  \[
    F_x
    \pause
    = \D{}{x}\big(y^2-x^2-\sin(xy)\big)
    \pause
    = 0 - 2x - y\cos(xy)
    \pause
    = - 2x - y\cos(xy)
  \]
  \pause
  \[
    F_y
    \pause
    = \D{}{y}\big(y^2-x^2-\sin(xy)\big)
    \pause
    = 2y - 0 - x\cos(xy)
    \pause
    = 2y - x\cos(xy)
  \]
  \pause
  \[
    \frac{dy}{dx}
    \pause
    = -\frac{F_x}{F_y}
    \pause
    = -\frac{- 2x - y\cos(xy)}{2y - x\cos(xy)}
    \pause
    = \frac{2x + y\cos(xy)}{2y - x\cos(xy)}
  \]
\end{frame}

%------------------------------------------------------------------------------%
\begin{frame}{Derivação Implícita -- Três Variáveis}
  Suponha que $F(x, y, z)$ seja diferenciável e que a equação
  \[
    F(x, y, z) = 0
  \]
  defina $z$ como função de $(x, y)$, isto é, $z=f(x, y)$,
  determine
  \(\D{z}{x}\) e
  \(\D{z}{y}\)
\end{frame}

%------------------------------------------------------------------------------%
\begin{frame}{Derivação Implícita -- Justificativa}
\begin{columns}
  \column{0.5\linewidth}
  \begin{align*}
    \onslide<+->{F\big(x, y, z(x, y)\big)                               & = 0 \\[2mm]}
    \onslide<+->{\D{} {x\vphantom{y}}\left[F\big(x, y, z(x, y)\big)\right]          & = 0 \\[2mm]}
    \onslide<+->{\D{F}{x\vphantom{y}}\D{x}{x} + \D{F}{y}\D{y}{x} + \D{F}{z}\D{z}{x} & = 0 \\[2mm]}
    \onslide<+->{\D{F}{x\vphantom{y}} + 0 + \D{F}{z}\D{z}{x}                        & = 0}
  \end{align*}
  \column{0.5\linewidth}
  \begin{align*}
    \onslide<+->{F_x + F_z \D{z}{x\vphantom{y}} & = 0                   \\[2mm]}
    \onslide<+->{F_z \D{z}{x\vphantom{y}}       & = - F_x               \\[2mm]}
    \onslide<+->{\D{z}{x\vphantom{y}}           & = - \frac{F_x}{F_z} \qquad \text{desde que } F_z \neq 0}
  \end{align*}
\end{columns}
\end{frame}

%------------------------------------------------------------------------------%
\begin{frame}{Derivação Implícita -- Justificativa}
  \begin{columns}
    \column{0.5\linewidth}
    \begin{align*}
      F\big(x, y, z(x, y)\big)                               & = 0 \\[2mm]
      \D{}{y}\left[F\big(x, y, z(x, y)\big)\right]           & = 0 \\[2mm]
      \D{F}{x}\D{x}{y} + \D{F}{y}\D{y}{y} + \D{F}{z}\D{z}{y} & = 0 \\[2mm]
      0 + \D{F}{y} + \D{F}{z}\D{z}{y}                        & = 0
    \end{align*}
    \column{0.5\linewidth}
    \begin{align*}
      F_y + F_z \D{z}{y} & = 0                   \\[2mm]
      F_z \D{z}{y}       & = - F_y               \\[2mm]
      \D{z}{y}           & = - \frac{F_y}{F_z} \qquad \text{desde que } F_z \neq 0
    \end{align*}
  \end{columns}
\end{frame}

%------------------------------------------------------------------------------%
\begin{frame}{Derivação Implícita -- Três Variáveis -- Fórmulas}
  Suponha que $F(x, y, z)$ seja diferenciável e que a equação
  \[
    F(x, y, z) = 0
  \]
  defina $z$ como função de $(x, y)$, então
  \[
    \D{z}{x} = -\frac{F_x}{F_z}
  \]
  \[
    \D{z}{y} = -\frac{F_y}{F_z}
  \]
  desde que \(F_z \neq 0\)
\end{frame}

%------------------------------------------------------------------------------%
\addtocounter{exemplo}{1}
\begin{frame}{Exemplo \theexemplo}
  Calcule \; \(\D{z}{x}\) e \(\D{z}{y}\) \; se \; \(x^3 + z^2 + ye^{xz} + z\cos(y) = 0\)
\end{frame}

\begin{frame}{Exemplo \theexemplo}
  \[
    F(x, y, z) = x^3 + z^2 + ye^{xz} + z\cos(y)
  \]
  \pause
  \[
    F_x
    \pause
    = \D{}{x}\left[x^3 + z^2 + ye^{xz} + z\cos(y)\right]
    \pause
    = 3x^2 + 0 + ye^{xz}z + 0
    \pause
    = 3x^2 + yze^{xz}
  \]
  \pause
  \[
    F_y
    \pause
    = \D{}{y}\left[x^3 + z^2 + ye^{xz} + z\cos(y)\right]
    \pause
    = 0 + 0 + e^{xz} -z\sin(y)
    \pause
    = e^{xz} -z\sin(y)
  \]
  \pause
  \[
    F_z
    \pause
    = \D{}{z}\left[x^3 + z^2 + ye^{xz} + z\cos(y)\right]
    \pause
    = 0 + 2z + ye^{xz}x + \cos(y)
    \pause
    = 2z + yxe^{xz} + \cos(y)
  \]
\end{frame}

\begin{frame}{Exemplo \theexemplo}
  \[
    \D{z}{x}
    = -\frac{F_x}{F_z}
    \pause
    = -\frac{3x^2 + yze^{xz}}{2z + yxe^{xy} + \cos(y)}
  \]
  \pause
  \[
    \D{z}{y}
    = -\frac{F_y}{F_z}
    \pause
    = -\frac{e^{xz} -z\sin(y)}{2z + yxe^{xy} + \cos(y)}
  \]
\end{frame}

%------------------------------------------------------------------------------%
\section{Exemplos}

%------------------------------------------------------------------------------%
\addtocounter{exemplo}{1}
\begin{frame}{Exemplo \theexemplo}
  Considerando que a equação
  \[
    x^3 + y^3 + z^3 = 3xyz + 1
  \]
  defina $z$ como função de $x$ e $y$. Encontre, se possível,
  \(\D{z}{x}\) e \(\D{z}{y}\)
  nos pontos \((1, 1, 1)\) e \((1, 3, 2)\)

  \pause
  \vspace{\baselineskip}
  Note que
  \(F(x,y,z) = x^3 + y^3 + z^3 - 3xyz - 1\)
\end{frame}

\begin{frame}{Exemplo \theexemplo}
  Sabemos que, se \(F_z \neq 0\)
  \[
    \D{z}{x} = -\frac{F_x}{F_z}
  \]
  \[
    \D{z}{y} = -\frac{F_y}{F_z}
  \]
\end{frame}

\begin{frame}{Exemplo \theexemplo}
  Calculando as derivadas parciais de $F(z, y, z)$

  \[
    \pause
    F_x
    \pause
    \;=\; \D{}{x}\left(x^3 + y^3 + z^3 - 3xyz - 1\right)
    \pause
    \;=\; 3x^2 - 3yz
  \]

  \[
    \pause
    F_y
    \pause
    \;=\; \D{}{y}\left(x^3 + y^3 + z^3 - 3xyz - 1\right)
    \pause
    \;=\; 3y^2 - 3xz
  \]

  \[
    \pause
    F_z
    \pause
    \;=\; \D{}{z}\left(x^3 + y^3 + z^3 - 3xyz - 1\right)
    \pause
    \;=\; 3z^2 - 3xy
  \]
\end{frame}

\begin{frame}{Exemplo \theexemplo}
  Avaliando as derivadas parciais no ponto \((1, 1, 1)\)

  \[
    F_z(1, 1, 1)
    \pause
    \;=\; \left(3z^2 - 3xy\right)\evalat{(1, 1, 1)}{}
    \pause
    \;=\;  3 - 3
    \pause
    \;=\; 0
  \]

  \pause
  \emph{Não é possível} calcular as derivadas neste ponto
\end{frame}

\begin{frame}{Exemplo \theexemplo}
  Avaliando as derivadas parciais no ponto \((1, 3, 2)\)

  \[
    \pause
    F_z(1, 1, 2)
    \pause
    \;=\; \left(3z^2 - 3xy\right)\evalat{(1, 3, 2)}{}
    \pause
    \;=\; 3\times 2^2 - 3 \times 1 \times 3
    \pause
    \;=\; 12 - 9
    \pause
    \;=\; 3
    \pause
    \;\neq\; 0
  \]

  \pause
  \emph{É possível} calcular as derivadas neste ponto

  \[
    \pause
    F_x(1, 1, 2)
    \pause
    \;=\; \left(3x^2 - 3yz\right)\evalat{(1, 3, 2)}{}
    \pause
    \;=\; 3 \times 1^2 - 3 \times 3 \times 2
    \pause
    \;=\; 3 - 18
    \pause
    \;=\; -15
  \]

  \[
    \pause
    F_y(1, 1, 2)
    \pause
    \;=\; \left(3y^2 - 3xz\right)\evalat{(1, 3, 2)}{}
    \pause
    \;=\; 3 \times 3^3 - 3 \times1 \times 2
    \pause
    \;=\; 27 - 6
    \pause
    \;=\; 21
  \]
\end{frame}

\begin{frame}{Exemplo \theexemplo}
  \[
    \D{z}{x}(1, 3, 2)
    \pause
    \;=\; -\frac{F_x(1, 3, 2)}{F_z(1, 3, 2)}
    \pause
    \;=\; -\frac{-15}{3}
    \pause
    \;=\; 5
  \]

  \[
    \pause
    \D{z}{y}(1, 1, 2)
    \pause
    \;=\; -\frac{F_y(1, 3, 2)}{F_z(1, 3, 2)}
    \pause
    \;=\; -\frac{21}{3}
    \pause
    \;=\; -7
  \]
\end{frame}

%------------------------------------------------------------------------------%
\addtocounter{exemplo}{1}
\begin{frame}{Exemplo \theexemplo}
  Dada a equação \; \(x^2y + y^2z + z^2x = 1\) \;
  encontre \; \(\D{z}{x}\)
\end{frame}

\begin{frame}{Exemplo \theexemplo}
  \onslide<+->{Derivando ambos os lados por $x$, tratando $z$ como função \(z=z(x, y)\)}
\begin{columns}
\column{0.6\linewidth}
  \begin{align*}
  \onslide<+->{x^2y + y^2z + z^2x                            & = 1          \\[2mm]}
  \onslide<+->{\D{}{x}\left(x^2y + y^2z + z^2x\right)        & = \D{1}{x}   \\[2mm]}
  \onslide<+->{\D{}{x}(x^2y) + \D{}{x}(y^2z) + \D{}{x}(z^2x) & = 0          \\[2mm]}
  \onslide<+->{2xy + y^2\D{z}{x} + 2zx\D{z}{x} + z^2         & = 0          }
\end{align*}
\column{0.4\linewidth}
  \begin{align*}
  \onslide<+->{\D{z}{x}\left(y^2 + 2zx\right) & = -2xy - z^2           \\[3mm]}
  \onslide<+->{\D{z}{x}                       & = -\frac{2xy + z^2}{y^2 + 2zx}}
\end{align*}
\end{columns}
\end{frame}

%------------------------------------------------------------------------------%
\section{Lista Mínima}

%------------------------------------------------------------------------------%
\begin{frame}{Lista Mínima}
  Cálculo Vol. 2 do Thomas 12\textsuperscript{a} ed. -- Seção 14.4
  \begin{enumerate}
    \item Estudar o texto da seção
    \item Resolver os exercícios: 26, 28, 30, 32
  \end{enumerate}

  \vfill
  Atenção: A prova é baseada no livro, não nas apresentações
\end{frame}

%------------------------------------------------------------------------------%
\end{document}
%------------------------------------------------------------------------------%
